\subsubsection{理論推導}

在波耳模型中，電子繞行原子核做圓周運動，其向心力由庫倫靜電力提供：

\[
\frac{mv^2}{r}=\frac{1}{4\pi\varepsilon_0}\frac{e^2}{r^2}
\]

其中 $m$ 為電子質量，$v$ 為電子速度，$r$ 為軌道半徑。

\vspace{0.3cm}

波耳提出角動量量子化假設：

\[
mvr = n\hbar \qquad (n=1,2,3,\dots)
\]

其中 $\hbar=\frac{h}{2\pi}$。

\vspace{0.3cm}

聯立兩式可求得電子允許的軌道半徑：

\[
r_n = \frac{4\pi\varepsilon_0 n^2 \hbar^2}{m e^2}
\]

當 $n=1$ 時得到基態半徑：

\[
a_0 = 0.529 \text{ Å}
\]

稱為\textbf{波耳半徑}。

\vspace{0.3cm}

電子的總能量為動能與位能之和：

\[
E = \frac{1}{2}mv^2 - \frac{1}{4\pi\varepsilon_0}\frac{e^2}{r}
\]

代入條件後可得氫原子能階：

\[
E_n = -\frac{13.6}{n^2}\ \text{eV}
\]

\vspace{0.3cm}

當電子由高能階 $n_i$ 躍遷至低能階 $n_f$ 時，會放出光子：

\[
\Delta E = E_i - E_f = h\nu
\]
換算為波長可得到氫原子的\textbf{芮得柏公式}：

\[
\frac{1}{\lambda} = R_H
\left(
\frac{1}{n_f^2} - \frac{1}{n_i^2}
\right)
\]

其中 $R_H = 1.097\times10^{-2}\, nm^{-1}$。